\chapter{Veille réglementaire et documentaire}
\section{Ma méthode de vérification}

Pour préparer la séance, je dois distinguer les règles juridiques, les statistiques d’accidentalité et les recommandations de prévention. Je privilégie donc les sources institutionnelles et je vérifie la date des informations avant de les utiliser dans le diaporama.

Mes principales sources sont :

— Légifrance et le Code de la route pour les textes applicables, notamment l’article R.412-6-1

sur le téléphone ;

— l’ONISR pour les bilans d’accidentalité et les observations de comportements ;

— Sécurité routière et Service-Public.fr pour les informations accessibles sur les règles et

les sanctions ;

— l’INRS, la CNAM et la MSA pour le risque routier professionnel et l’organisation des

déplacements en mission.

Cette veille a une utilité directe : elle me permet de corriger ou d’actualiser une donnée avant de la présenter. Elle m’aide aussi à choisir des informations utiles au public plutôt qu’à multiplier les chiffres.

Par exemple, pour les distracteurs, je dois faire attention aux termes utilisés. L’inattention, l’usage du téléphone et les distracteurs technologiques ne sont pas toujours regroupés de la même manière dans les statistiques. Je dois donc présenter chaque chiffre avec son année et son périmètre, sans lui faire dire plus que la source. De la même façon, les données concernant les véhicules utilitaires ne doivent pas être présentées comme si tous les accidents étaient liés au téléphone.

Cette rigueur fait partie de ma posture de futur ECSR. Si une information évolue entre la préparation du document écrit et la date de l’animation, c’est la donnée la plus récente et la mieux sourcée qui doit être utilisée.

[À compléter : Avant la remise définitive : ajouter un tableau court avec les documents réellement consultés, leur date, le point vérifié et l’incidence sur la séance.]
